\section*{Executive Summary}
\phantomsection\addcontentsline{toc}{section}{Executive Summary}
\label{sec:exec}

% ---- Condensed rework of your draft (your version: 01_executive_summary_draft.tex).
% ---- Green numbers are from the results pipeline tables.

This project designed, built and field-tested a controller for the `Power Shuffler', and showed on a real solar string that it recovers power a shaded string would otherwise lose. Shading one panel in a series string limits the current of the whole string, and the bypass diodes used to protect it simply discard the shaded panel's power. The shuffler is a serial partial power converter (S-PPC), a type of differential power processor, that instead moves the difference in power between series-connected modules to balance their voltages. A sister report covers the shuffler hardware; this report covers its control system, from the analytical model through to closed-loop validation on the bench and in the field.

Each shuffler is an interleaved buck-boost converter built from decoupling capacitors, a coupled inductor and two half-bridges. Each shuffler shares PV modules with its neighbours, so their dynamics are coupled and cannot be described by a set of single-input single-output transfer functions. The controller is therefore built on a multiple-input multiple-output state-space model with three states per shuffler (module-voltage difference, shuffled current and circulating current) and two inputs: common-mode duty (set by the two half-bridge duties together, which shuffles power between modules) and differential duty (their difference, which controls the circulating current). The development process followed a unique spin on a systems engineering approach. The analytical model set the controller architecture. A Simulink model of a four-module string validated these control loops, and its requirements drove the choice of sensors and microcontroller. The same C control core was then implemented inside a hardware-accurate QSPICE model, adding sensor noise and sampling delay, before going into the firmware unchanged. The final product is a single board with two shufflers, Hall-effect current sensors, power-monitor voltage sensors and an STM32G474RE microcontroller, housed in a 3D-printed enclosure with a Raspberry Pi 4B that streams telemetry to a timeseries database.

On the bench (NGE103 supply, DL3021 electronic load) the closed loop held the module-voltage differences within $\pm0.010\V$ and the circulating current within $\pm0.07\A$ across 22 settled test segments at 200--400\,kHz. The 44\,Hz telemetry is too slow to resolve the step response on hardware, so it is taken from the simulation models. For a 1\,V step at 30\,\% gain and 200\,kHz, QSPICE gives a rise time of 6.5\ms{} and a 2\,\% settling time of 11.9\ms{} with no overshoot, corroborated by the analytical model (7.6\ms{} and 13.6\ms). Of the switching frequencies tested, 300\,kHz was the best compromise; 400\,kHz gave marginal gains at higher losses with almost no change in response from gain increases.

The most significant result is from the solar test: three 80\,W panels without bypass diodes in series into a 61\,V constant-voltage load. With the top module shaded to almost zero current (a shirt over about a third of the panel), a plain series string would have delivered no power at all. With the shufflers running, the string delivered up to 99\W{} during the shading event and 10.6\Wh{} over 36 minutes of testing, nearly all of it energy that would otherwise have been lost.

The main limitation is scope: one board with two shufflers that are hard-wired together, so they cannot be driven independently and boards cannot be stacked, which leaves scalability untested. The top-level INA239 (U10) does not work, so $V_{M1}$ was inferred from the electronic load; the gains have to be re-derived in simulation whenever the switching frequency changes; and poor weather on the test day meant the closed loop was only partly exercised against real panel dynamics. The recommended next steps are a stackable board with independently driven shufflers, to test strings of more than three modules, and a study of other applications for S-PPCs that may offer more impact to industry than solar.
